%% LyX 2.5.1 created this file.  For more info, see https://www.lyx.org/.
%% Do not edit unless you really know what you are doing.
\documentclass[12pt]{article}
\usepackage[T1]{fontenc}
\usepackage{enumitem}
\usepackage{amsmath}
\usepackage{amssymb}
\usepackage[a4paper]{geometry}
\geometry{verbose,headheight=14pt,footskip=1cm}
\usepackage{graphicx}
\usepackage[english]{babel}

\makeatletter
%%%%%%%%%%%%%%%%%%%%%%%%%%%%%% Textclass specific LaTeX commands.
\newlength{\lyxlabelwidth}      % auxiliary length 

%%%%%%%%%%%%%%%%%%%%%%%%%%%%%% User specified LaTeX commands.
% Better than enumerate, more flexible
%\usepackage{enumerate}
\newcommand{\encabezadoexamen}{
    \begin{center}
        \textbf{\Large Name of the Institution} \\
        \vspace{0.5cm}
        \textbf{\Large Name of the Subject} \\
        \vspace{0.3cm}
        \textbf{\Large PARTIAL / FINAL EXAM} \\
        \vspace{0.2cm}
        \rule{\textwidth}{0.8pt}
        \end{center}
        \vspace{0.5cm}
        \begin{flushleft}
        \textbf{Student Name:} \hrulefill \\
        \vspace{0.2cm}
        \textbf{Date:} \hrulefill \quad \textbf{Group:} \hrulefill \\
        \vspace{0.2cm}
        \textbf{Total Score:} \hrulefill \quad \textbf{Score Obtained:} \hrulefill \\
        \end{flushleft}
        \vspace{0.5cm}
        \rule{\textwidth}{0.4pt}
        \vspace{0.5cm}
    }
% Solution: use \item[] for the introductory text (without label)
\newenvironment{pregunta}[1]{%
    \vspace{0.5cm}%
    \noindent\textbf{Question #1}%
    \par%
    \begin{enumerate}[label=\alph*)]%
}{%
    \end{enumerate}%
    \vspace{0.5cm}%
}

\makeatother

\begin{document}
\encabezadoexamen

\section*{General Instructions}
\begin{itemize}[label=$\bullet$]
\item Read all questions carefully before answering. 
\item Use blue or black pen. Pencil answers are not accepted. 
\item The duration of the exam is \textbf{90 minutes}. 
\item The score for each question is indicated at the beginning of the question. 
\item Clarity, order, and justification of the answers will be valued. 
\end{itemize}
\vspace{1cm}

\begin{pregunta}{1 (2 points)} 

\item[] Answer the following questions on topic X.
% Use \item[] for unnumbered text

\item Explain briefly the concept of \textit{Fundamental Theorem of
Calculus}. 

\item Mention two practical applications of derivatives in
engineering. \end{pregunta}

\begin{pregunta}{2 (3 points)} 

\item[] Solve the following linear algebra problem. 

\item Calculate the inverse of the matrix $A=\begin{pmatrix}1 & 2\\
3 & 4
\end{pmatrix}$. 

\item Determine whether the vectors $v_{1}=(1,0,1)$, $v_{2}=(0,1,1)$
and $v_{3}=(1,1,0)$ are linearly independent. \end{pregunta}

\begin{pregunta}{3 (2.5 points)} 

\item[] Prove the following proposition. 

\item Prove that if a function $f(x)$ is continuous on a closed
interval $[a,b]$ and differentiable on the open interval $(a,b)$, then
there exists at least one point $c\in(a,b)$ such that $f'(c)=\frac{f(b)-f(a)}{b-a}$.
(Mean Value Theorem) \end{pregunta}

\begin{pregunta}{4 (2.5 points)} 

\item[] Application problem. 

\item A company produces $Q$ units of a product, where the total
cost of production is given by the function $C(Q)=0.01Q^{2}+5Q+100$.
If the selling price per unit is \$20, determine the quantity
of units the company should produce to maximize its profit.
\end{pregunta}

\vspace{2cm}

\begin{center}
\textit{Good luck!} 
\par\end{center}
\end{document}
